\documentclass[5pt, a4paper]{article}

\usepackage[utf8]{inputenc}
\usepackage[T1]{fontenc}
\usepackage{textcomp}
\usepackage{amsmath, amssymb}
\usepackage{multicol}
\usepackage{proof}
\usepackage{enumitem}
\usepackage{fancyhdr}
\usepackage[margin=1in]{geometry}
% figure support
\usepackage{import}
\usepackage{xifthen}
\pdfminorversion=7

\usepackage[most]{tcolorbox}
\usepackage{pdfpages}
\usepackage{transparent}
\newcommand{\incfig}[1]{%
    \def\svgwidth{\columnwidth}
    \import{./figures/}{#1.pdf_tex}
}

\tcbset{
colback=gray!5,
colframe=blue!50!white,
coltitle=black,
arc=0pt,
boxrule=-0.1mm,
fonttitle=\bfseries,
coltitle=black,
left=2.5mm,
leftrule=1mm,
right=3.5mm,
colbacktitle=black!10!white,
toptitle=0.75mm,
bottomtitle=0.5mm,
}
\title{ST2131 Probability}
\date{Y2S2 2026}
\author{Nawwaf Sudi}
\begin{document}
    \maketitle
    \newpage
    \tableofcontents
    \newpage
    \section{Combinatorial Analysis and Probability}
    Let's begin with an intuitive introduction. Suppose you have a task, that
    takes a couple of steps to complete. Suppose for the $i$-th task, we have
    $m_i$ possible ways to perform it. When determining the total amount of ways
    you have to complete the task, we follow the multiplication principle. 
    \begin{tcolorbox}[title=Multiplication Principle]
        For an experiment consisting of $n$ different independent events, where
        $m_1,m_2,\ldots,m_n$ are the number of possible outcomes of the event,
        the total number of possible outcomes $N$ is given by \[
            N = m_1m_2\ldots m_n
        .\] 
    \end{tcolorbox}

    Suppose, again, that we are to pick a ball out of $n$ urns, where each urns
    contain $m_n$ number of balls. The total amount of possible balls that we
    could pick is given by the addition principle.
    \begin{tcolorbox}[title=Addition Principle]
        For $n$ sets of outcomes that each has $m_i$ elements, where $i$ is the
        $i$-th set, and all the set is mutually exclusive, the total number of
        outcomes $N$ is  \[
            N = m_1+m_2+\ldots + m_n
        .\] 
    \end{tcolorbox}

    \section{Discrete Random Variable}

    \section{Continuous Random Variable}
    \begin{tcolorbox}[title=Continuous Random Variable]
        A random variable $X$ is said to be continuous if there exists a
        nonnegative function $f_X$ defined for all $x \in \mathbb{R}$ such that
        for any set $B \subseteq \mathbb{R}$, \[
            P(X \in B) = \int_B f_X(x) dx
        .\] 
        $f_X$ is refered to as the Probability Desity Function (PDF). Note that
         $X$ takes on an uncountably large number of support.
    \end{tcolorbox}
    \paragraph{Example} For a continuous random variable $X$, and a real
    interval $B=[a,b]$, the probability that $X$ takes on the value inside $B$
    is given by  \[
            P(a \le X \le b) = \int_a^b f_X(x) dx
    .\] 
    
    \begin{tcolorbox}[colframe=blue!20!white]
        Another way of understanding the PDF of a continuous random variable is
        thus: Suppose a very small interval $\epsilon$ and a PDF $f_X$ continous
        at a particular point  $x$. The probability that $X$ 
        takes on the value $x - \frac{\epsilon}{2} \le X \le x +
        \frac{\epsilon}{2}$ is \[
        P(x - \frac{\epsilon}{2} \le x \le x + \frac{\epsilon}{2})
        = \int_{x-\frac{\epsilon}{2}}^{x+\frac{\epsilon}{2}} f_X(x) dx
        .\] 
        Since $\epsilon$ is small, then we can model the probability that $X$ 
        falls into the interval as $\epsilon f_X(x)$, and as such the function
        can be thought of as the probability that the random variable would
        take on values near $x$.

        Also, notice that if $\epsilon \to 0$ (ie. $a,b \to x$), $X$ would take
        on the value of $x$, and \[
        P(X = x)
        = \int_{x-\frac{\epsilon}{2}}^{x+\frac{\epsilon}{2}} f_X(x) dx = 0
        .\] 
        We can conclude that the probability that a continuous random variable
        takes on a single particular value is zero
    \end{tcolorbox}

    \begin{tcolorbox}[title= Continuous Distribution Function]
        We define the CDF for a continuous random variable the exact same way
        that we define it for the discreete random variable, which is \[
        F_X(x) = P(X \le x)
        .\] 
    \end{tcolorbox}
    \begin{tcolorbox}[title=Properties of CDF]
        \begin{enumerate}
            \item $F_X(x)$ is a nondecreasing function.
            \item $\lim_{b \to \infty}F_X(b) =1$
            \item $\lim_{b \to -\infty}F_X(b) =0$ 
            \item $F_X(x)$ is right continuous ($\lim_{x \to b^+} F_X(x) =
                F_X(b)$)
        \end{enumerate}
    \end{tcolorbox}

    Using the above properties, and the definition of a continuous random
    variables, we can write the following colloraries:

    \[P(a < X \le b) = F_X(b) - F_X(a)\]

    \paragraph{Proof}

    We write $\{X\le b\} = \{X \le  a\} \cup \{a < X < b\}$. Thus, 
    \begin{align*}
        P(X \le  b) &= P(X \le  a) + P(a < X \le b) \\
        P(a < X \le b) &= P(X \le b) - P(X \le b)
    .\end{align*}

    \[P(X <b) = \lim_{n \to \infty} \left( b - \frac{1}{n} \right) \]

    \paragraph{Proof}

        \[P(X=a) = F_X(a) - P(X < a).\] 

    \paragraph{Proof} aa


    \begin{tcolorbox}[title= q-th Quantile]
        A $q$-th quantile of a random variable $X$ is defined as the number
        $z_q$ such that  \[
        P(X \le z_q) = q
        ,\] with $0\le q\le 1$  
    \end{tcolorbox}
    \begin{tcolorbox}[title=Expectation of a Continuous Random Variable]
        The expectation of a continuous random variable $X$ is given by  \[
            E(X) = \int_{-\infty}^{\infty} xf_X(x) dx
        .\] 
        The variance of $X$ is given by  \[
            var(X) = \int_{-\infty}^{\infty}\left(x - E(X)\right)^2 f(x)dx
        .\] 
    \end{tcolorbox}
    \begin{tcolorbox}[title=Tail-Sum Formula]
        The expected value of a continuous random variable is given by \[
            E(X) = \int_{0}^{\infty}P(X\ge x)dx = \int_{0}^{\infty}P(X>x)dx
        .\] 
    \end{tcolorbox}

    \begin{tcolorbox}[title=Normal Distribution]
        A random variable $X$ is sad to be normally distributed with parameters
         $\mu$ and $\sigma^2$ if it's probability density function is given by
          \[
              f_X(x) = \frac{1}{\sqrt{2\pi}\sigma
              }e^{-\frac{(x-\mu)^2}{2\sigma^2}}
         .\]
         Denoted by N($\mu,\sigma^2$).
         N(0,1) is called a standard normal, denoted by $Z ~N(0,1)$. It's PDF
         and CDF is given by  \[
             \phi(X) = \frac{1}{\sqrt{2\pi}e^{-\frac{x^2}{2}} }
         .\] 
         \[
             \Phi(x) = \frac{1}{\sqrt{2\pi} } \int_{-\infty}^{x}
             e^{-\frac{y^2}{2}}dy
         .\] 
    \end{tcolorbox}
    \begin{tcolorbox}[title=Normal Distribution Property, colframe=blue!20!white]
        If $Y ~ N(\mu,\sigma^2)$, then $E(Y) = \mu$ and $var(Y) = \sigma^2$. We
        can write the relation between $Y$ and $Z$ like so: \[
        P(a \le Y\le b) = P(\frac{a-\mu}{\sigma} \le Z \le \frac{b -
        \mu}{\sigma}) = \Phi(\frac{b-\mu}{\sigma}) - \Phi(\frac{a-\mu}{\sigma})
        .\] 
    \end{tcolorbox}
    \begin{tcolorbox}[title=Standard Normal Distribution Properties,
        colframe=blue!20!white]
        \begin{enumerate}
            \item $P(Z \ge 0) = P(Z \le 0) = \frac{1}{2}$
            \item $-Z ~ N(0,1)$ 
            \item $P(Z\le x) = 1 - P(Z > x)$
            \item $P(Z \le -x) = P(Z > x)$
            \item If $Y~N(\mu,\sigma^2)$, then $X = \frac{x- \mu}{\sigma^2} ~
                N(0,1)$
            \item If $X ~ N(0,1)$, then $Y=aX + b ~ N(b,a^2)$
        \end{enumerate}
        
    \end{tcolorbox}

    \begin{tcolorbox}[title=Exponential Distribution]
        A random variable is said to be exponentially distributed if the PDF is
        given by \[
        f_X(x) = \begin{cases}
            \lambda e^{-\lambda x} & x \ge 0 \\
            0 & x < 0
        \end{cases}
        .\] 
        Denoted by Exp($\lambda$). It's CDF is given by \[
        f_X(x) = \begin{cases}
            0 & x \le 0 \\
            1 - e^{-\lambda x} & x > 0 \\
        \end{cases}
        .\]  The expected value and variance is given by \[
        E(X) = \frac{1}{\lambda} \qquad var(X) = \frac{1}{\lambda^2}
        .\] 

    \end{tcolorbox}
\end{document}
